\par\Needspace{6\baselineskip}\medskip
\begin{center}\hypertarget{native-245-article174}{}\pdfbookmark[2]{174.}{article-174}174.\end{center}
Si corporis coelestis orbita secundum methodos in Sectionibus praec. traditas
e tribus quatuorve positionibus geocentricis talibus determinata est, quae
ipsae singulae ad normam art.\ praec. e compluribus observationibus petitae
fuerant. orbita ista inter omnes hasce observationes medium quasi tenebit,
neque in differentiis inter locos observatos et calculatos ullum ordinis
vestigium remanebit, quod per elementorum correctionem tollere vel sensibiliter
extenuare liceret. Iam quoties tota observationum copia intervallum temporis
non nimis magnum complectitur, hoc modo consensum exoptatissimum elementorum
cum omnibus observationibus assequi licebit, si modo tres quatuorve positiones
quasi normales scite eligantur. In determinandis orbitis cometarum
planetarumve novorum, quorum observationes annum unum nondum egrediuntur,
ista ratione plerumque tantum proficiemus, quantum ipsa rei natura permittit.
Quoties itaque orbita determinanda angulo considerabili versus eclipticam
inclinata est, in genere tribus observationibus superstruetur, quas quam
remotissimas ab invicem eligemus: si vero hoc pacto in aliquem casuum supra
exclusorum (artt.\ 160---162) fortuito incideremus, aut quoties orbitae
inclinatio nimis parva videtur, determinationem ex positionibus quatuor
praeferemus, quas itidem quam remotissimas ab invicem accipiemus.

Quando autem iam adest observationum series longior plures annos complectens,
plures inde positiones normales derivari poterunt: quamobrem praecisioni
maximae male consuleremus, si ad orbitae determinationem tres tantum
quatuorve positiones excerperemus, omnesque reliquas omnino negligeremus.
Quin potius in tali casu, si summam praecisionem assequi propositum est,
operam dabimus, ut positiones exquisitas quam plurimas congeramus, atque in
usum vocemus. Tunc itaque aderunt data plura, quam ad incognitarum
determinationem requiruntur: sed omnia ista data erroribus utut exiguis
obnoxia erunt, ita ut generaliter impossibile sit, omnibus ex asse satisfacere.
Iam quum nulla adsit ratio, cur ex hisce datis sex haec vel illa tamquam
absolute exacta consideremus, sed potius, secundum probabilitatis principia,
in \SourceLeaf{246}cunctis promiscue errores maiores vel minores aeque possibiles
supponere oporteat; porro quum generaliter loquendo errores leviores saepius
committantur quam graviores: manifestum est, orbitam talem, quae dum sex
datis ad amussim satisfacit a reliquis plus minusve deviat, principiis calculi
probabilitatis minus consentaneam censendam esse, quam aliam, quae, dum ab
illis quoque sex datis aliquantulum discrepat, consensum tanto meliorem cum
reliquis praestat. Investigatio orbitae sensu stricto \emph{maximam} probabilitatem
prae se ferentis a cognitione legis pendebit, secundum quam errorum
crescentium probabilitas decrescit: illa vero a tot considerationibus vagis vel
dubiis --- physiologicis quoque --- pendet, quae calculo subiici nequeunt, ut
huiusmodi legem vix ac ne vix quidem in ullo astronomiae practicae casu rite
assignare liceat. Nihilominus indagatio nexus inter hanc legem orbitamque
maxime probabilem, quam summa iam generalitate suscipiemus, neutiquam pro
speculatione sterili habenda erit.
